%% GENERATED from PAPER_SUBMISSION.md (round-7) -- hand-fixed conversion; see BUILD_NOTES.md
\section{Ethics and Data Governance}\label{sec:10}

This work uses only the public, pseudonymized \textbf{Amazon Reviews 2023} dataset (Hou et al., 2024;
This work uses only the public, pseudonymized \textbf{Amazon Reviews 2023} dataset (Hou et al., 2024;
McAuley Lab), obtained from its official distribution and publicly released for research by the McAuley Lab --- whose maintainers state they are not in a position to assign a license or dictate usage terms --- and we redistribute derived split CSVs on that basis with attribution, removing them on maintainer request. Scope of
data actually consumed: interaction tuples (pseudonymous user identifiers --- remapped to dense
integers by our preprocessing --- item identifiers, and timestamps) and \textbf{item metadata text}
(titles, categories, and store/seller name --- the AR2023 \texttt{store} field; an earlier label said ``brand'', corrected 2026-07-20: the frozen cache text template's literal \texttt{brand:} prefix is retained as-is inside the cached strings and documented as a misnomer, because regenerating the caches would invalidate every frozen-text result) for the frozen text encoders. We do \textbf{not} process review text
bodies, ratings-as-text, or product images, and we make no attempt at re-identification of any
user. Released artifacts follow the same boundary: per-user evaluation sidecars contain only
remapped integer IDs, target item IDs, and rank outcomes; the raw dataset is \textbf{not
redistributed} (our releases carry SHA256 hashes and regeneration scripts instead,
\texttt{RELEASE\_MANIFEST.json}). No new data was collected and no interaction with human subjects took
place, so no institutional review determination was sought or is claimed; the data are public, platform-pseudonymized product reviews, and we state that as a fact rather than adjudicating any jurisdiction's requirements (per-user sidecars carry only the platform's hashed user identifiers, an item index, and a rank). Finally, this is an offline evaluation
study: recommender systems deployed on such data can amplify popularity and exposure biases ---
our long-tail analyses quantify one aspect of that concern --- and nothing in this paper
constitutes a deployment claim.
data actually consumed: interaction tuples (pseudonymous user identifiers --- remapped to dense
integers by our preprocessing --- item identifiers, and timestamps) and \textbf{item metadata text}
(titles, categories, and store/seller name --- the AR2023 \texttt{store} field; an earlier label said ``brand'', corrected 2026-07-20: the frozen cache text template's literal \texttt{brand:} prefix is retained as-is inside the cached strings and documented as a misnomer, because regenerating the caches would invalidate every frozen-text result) for the frozen text encoders. We do \textbf{not} process review text
bodies, ratings-as-text, or product images, and we make no attempt at re-identification of any
user. Released artifacts follow the same boundary: per-user evaluation sidecars contain only
remapped integer IDs, target item IDs, and rank outcomes; the raw dataset is \textbf{not
redistributed} (our releases carry SHA256 hashes and regeneration scripts instead,
\texttt{RELEASE\_MANIFEST.json}). No new data was collected and no interaction with human subjects took
This work uses only the public, pseudonymized \textbf{Amazon Reviews 2023} dataset (Hou et al., 2024;
McAuley Lab), obtained from its official distribution and publicly released for research by the McAuley Lab --- whose maintainers state they are not in a position to assign a license or dictate usage terms --- and we redistribute derived split CSVs on that basis with attribution, removing them on maintainer request. Scope of
data actually consumed: interaction tuples (pseudonymous user identifiers --- remapped to dense
integers by our preprocessing --- item identifiers, and timestamps) and \textbf{item metadata text}
(titles, categories, and store/seller name --- the AR2023 \texttt{store} field; an earlier label said ``brand'', corrected 2026-07-20: the frozen cache text template's literal \texttt{brand:} prefix is retained as-is inside the cached strings and documented as a misnomer, because regenerating the caches would invalidate every frozen-text result) for the frozen text encoders. We do \textbf{not} process review text
bodies, ratings-as-text, or product images, and we make no attempt at re-identification of any
user. Released artifacts follow the same boundary: per-user evaluation sidecars contain only
remapped integer IDs, target item IDs, and rank outcomes; the raw dataset is \textbf{not
redistributed} (our releases carry SHA256 hashes and regeneration scripts instead,
\texttt{RELEASE\_MANIFEST.json}). No new data was collected and no interaction with human subjects took
place, so no institutional review determination was sought or is claimed; the data are public, platform-pseudonymized product reviews, and we state that as a fact rather than adjudicating any jurisdiction's requirements (per-user sidecars carry only the platform's hashed user identifiers, an item index, and a rank). Finally, this is an offline evaluation
study: recommender systems deployed on such data can amplify popularity and exposure biases ---
our long-tail analyses quantify one aspect of that concern --- and nothing in this paper
constitutes a deployment claim.
study: recommender systems deployed on such data can amplify popularity and exposure biases ---
our long-tail analyses quantify one aspect of that concern --- and nothing in this paper
constitutes a deployment claim.
